% !TeX program = XeLaTeX
% !TeX root = ../pujavidhanam.tex
\chapt{श्री-अखण्डदीप-पूजा-कल्पः}

\centerline{\small{(मूलम्—श्रीव्रतरत्नाकरः, प्रथमो भागः)}}

\sect{अखण्डदीपारोपण-विधिः}

\uvacha{सूत उवाच}

\twolineshloka*
{अखण्डदीपारोपस्य विधिं वक्ष्ये मुनीश्वराः}
{शृण्वन्तु श्रद्धया यूयम् आश्वयुग्दर्शतः परम्}

\twolineshloka*
{प्रतिपत्तिथिमारभ्य मासमेकं निरन्तरम्}
{अखण्डदीपस्संस्थाप्यः श्रीविष्णोः परितुष्टये}

\twolineshloka*
{सौवर्णं वा राजतं वा कांस्यं वाऽपि स्वशक्तितः}
{कृत्वा पात्रं गोघृतेन पूरयेदब्जतन्तुभिः}

\twolineshloka*
{ततः कुर्यान्मासमध्ये विच्छेदो न भवेद्यथा}
{कल्पोक्तेन प्रकारेण पूजयेत् प्रत्यहं द्विजाः}

\twolineshloka*
{पूजान्ते च प्रतिदिनं दीपं दद्याद्द्विजन्मने}
{हैमं वा राजतं वाऽपि कांस्यं वाऽपि सदक्षिणम्}

\twolineshloka*
{भोजयेच्च द्विजं भक्त्या मासमेवं समाचरेत्}
{अशक्तावपि मासान्ते सह वा भोजयेद्द्विजान्}

\twolineshloka*
{विप्रानुज्ञां पुरा लब्ध्वा सम्पूज्य च गणाधिपम्}
{कुर्यात् सङ्कल्पमथ च वृद्धिश्राद्धं समाचरेत्}

\twolineshloka*
{पुण्याहं वाचयेत्तत्र दीपमारोपयेत्ततः}
{कलशं शङ्खमभ्यर्च्य पीठपूजां विधाय च}

\onelineshloka*
{भक्त्या समर्चयेद्विप्रो ज्योतिष्येव जनार्दनम्}

\input{purvanga/vighneshwara-puja}

\sect{प्रधान-पूजा — श्रीभूमिनीलासमेत-श्रीमहाविष्णु-पूजा}

\twolineshloka*
{शुक्लाम्बरधरं विष्णुं शशिवर्णं चतुर्भुजम्}
{प्रसन्नवदनं ध्यायेत् सर्वविघ्नोपशान्तये}

\dnsub{सङ्कल्पः}

ममोपात्त-समस्त-दुरित-क्षयद्वारा श्री-परमेश्वर-प्रीत्यर्थं शुभे शोभने मुहूर्ते अद्य ब्रह्मणः
द्वितीयपरार्धे श्वेतवराहकल्पे वैवस्वतमन्वन्तरे अष्टाविंशतितमे कलियुगे प्रथमे पादे
जम्बूद्वीपे भारतवर्षे भरतखण्डे मेरोः दक्षिणे पार्श्वे शकाब्दे अस्मिन् वर्तमाने व्यावहारिकाणां
प्रभवादीनां षष्ट्याः संवत्सराणां मध्ये \blank\see{app:samvatsara_names} नाम संवत्सरे \blank{} अयने
\blank{}-ऋतौ आश्विन-मासे कृष्णपक्षे प्रतिपदि शुभतिथौ \blank{} वासरयुक्तायाम्
\blank\see{app:nakshatra_names} नक्षत्र \blank\see{app:yoga_names} नाम योगे
\blank\see{app:karanam_names}-करण-युक्तायां च एवं गुण विशेषण विशिष्टायाम्
अस्याम् प्रतिपदि शुभतिथौ मम अज्ञाननिवृत्तिद्वारा सुज्ञानसिद्ध्यर्थं परलोके अन्धत्वनिवृत्त्यर्थं
श्रीभूमिनीलासमेत-श्रीमहाविष्णु-प्रीत्यर्थं अद्यप्रभृति आगामि-दर्शपर्यन्तं मासमेकं निरन्तरं
पुराणोक्तप्रकारेण अखण्डदीपे यथाशक्ति-ध्यानावाहनादिषोडशोपचारैः
श्रीभूमिनीलासमेत-श्रीमहाविष्णु-पूजां करिष्ये। तदङ्गं कलशपूजां च करिष्ये।

श्रीविघ्नेश्वराय नमः यथास्थानं प्रतिष्ठापयामि। शोभनार्थे क्षेमाय पुनरागमनाय च।\\
(गणपति-प्रसादं शिरसा गृहीत्वा)

\input{purvanga/aasana-puja}

\input{purvanga/ghanta-puja}

\begingroup
\renewcommand{\devaName}{देव}
\input{purvanga/kalasha-puja}
\endgroup

\input{purvanga/aatma-puja}

\input{purvanga/pitha-puja}

\input{purvanga/guru-dhyanam}

\dnsub{प्राण-प्रतिष्ठा}

असुनीते पुनरिति ऋचं पठित्वा तच्चक्षुर्देवहितम् इति मन्त्रेण देवस्य आज्येन नेत्रोन्मीलनं कृत्वा प्राणान् प्रतिष्ठापयामि इति प्राणप्रतिष्ठां कृत्वा दीपं प्रज्वाल्य वक्ष्यमाणप्रकारेण पूजयेत्॥

\sect{षोडशोपचार-पूजा}
\renewcommand{\devAya}{श्रीभूमिनीलासमेत-श्रीमहाविष्णवे नमः,}
\begin{center}

\twolineshloka*
{ध्यायेच्चतुर्भुजं देवं शङ्खचक्रगदाधरम्}
{श्रीभूमिनीलासहितं प्रसन्नमुखपङ्कजम्}
\textbf{अस्मिन् दीपे श्रीभूमिनीलासमेतं श्रीमहाविष्णुं ध्यायामि।}
\medskip

\twolineshloka*
{आवाहयामि लक्ष्मीशम् आदिमध्यान्तवर्जितम्}
{श्यामलं शान्तहृदयं ज्योतिष्यस्मिन् जनार्दनम्}
\textbf{श्रीभूमिनीलासमेतं श्रीमहाविष्णुम् आवाहयामि।}
\medskip

\twolineshloka*
{मुक्तावैडूर्यखचितं मुकुन्द मुनिसेवित}
{रत्नसिंहासनं चारु रमानाथ ददामि ते}
\textbf{\devAya{} आसनं समर्पयामि।}
\medskip

\twolineshloka*
{पाद्यं गृहाण भगवन् पङ्कजाक्ष परात्पर}
{हेमपात्रस्थितं शुद्धं हेतुत्रयविवर्जितम्}
\textbf{\devAya{} पाद्यं समर्पयामि।}
\medskip

\twolineshloka*
{अर्घ्यं गृहाण भगवन्न् अच्युतानन्त केशव}
{अक्षतादिसमायुक्तम् अर्घ्यं देहि मे पदम्}
\textbf{\devAya{} अर्घ्यं समर्पयामि।}
\medskip

\twolineshloka*
{अविशुद्धमिदं तोयं सुवर्णकलशस्थितम्}
{गृहाणाचमनार्थाय गोविन्द गरुडध्वज}
\textbf{\devAya{} आचमनीयं समर्पयामि।}
\medskip

\twolineshloka*
{मधुवैरिन् महादेव महनीयपदाम्बुज}
{मधुपर्कं गृहाणाद्य महापापविनाशन}
\textbf{\devAya{} मधुपर्कं समर्पयामि।}
\medskip

\twolineshloka*
{पञ्चामृतं गृहाणेदं पञ्चपातकनाशन}
{पञ्चाननसमाराध्य पञ्चबाणातिसुन्दर}
\textbf{\devAya{} पञ्चामृतस्नानं समर्पयामि।}
\medskip

\twolineshloka*
{तुङ्गा-गोदावरी-कृष्णा-गङ्गादिभ्यस्समाहृतम्}
{सलिलं विमलं देव कमलाधव गृह्यताम्}
\textbf{\devAya{} स्नानं समर्पयामि। स्नानानन्तरम् आचमनीयं समर्पयामि।}
\medskip

\twolineshloka*
{पीताम्बरयुगं देव दामोदर गुणार्णव}
{गृह्यतां करुणासिन्धो भक्त्या तुभ्यं समर्पितम्}
\textbf{\devAya{} वस्त्रयुग्मं समर्पयामि।}
\medskip

\twolineshloka*
{यज्ञोपवीतं परमं पवित्रं पापनाशनम्}
{दत्तं मया गृहाणेदं धारणार्थं दयानिधे}
\textbf{\devAya{} यज्ञोपवीतम् आभरणानि च समर्पयामि।}
\medskip

\twolineshloka*
{श्रीगन्धं कुङ्कुमोपेतं श्रीनिवास जगत्पते}
{विष्णो विलेपनार्थाय गृहाण गरुडध्वज}
\textbf{\devAya{} गन्धान् धारयामि।}
\medskip

\twolineshloka*
{अक्षतान् धवलान् देव शालिजान् शासितासुर}
{गृहाण कमलाकान्त देहि मे पदमक्षतम्}
\textbf{\devAya{} अक्षतान् समर्पयामि।}
\medskip

\twolineshloka*
{पङ्कजैः पारिजातैश्च मन्दारैर्मल्लिकासुमैः}
{तुलसीदलसंयुक्तैः पूजयामि पुरातनम्}
\textbf{\devAya{} पुष्पाणि समर्पयामि। पुष्पैः पूजयामि।}
\medskip

\end{center}

\dnsub{अङ्ग-पूजा}
\begin{longtable}{ll@{—}l}
१. & पावनाय नमः & पादौ पूजयामि\\
२. & गोविन्दाय नमः & गुल्फौ पूजयामि\\
३. & ज्योतिस्स्वरूपाय नमः & जङ्घे पूजयामि\\
४. & जगत्कुक्षये नमः & जानुनी पूजयामि\\
५. & उपेन्द्राय नमः & ऊरू पूजयामि\\
६. & कमलाक्षाय नमः & कटिं पूजयामि\\
७. & गुणनिधये नमः & गुह्यं पूजयामि\\
८. & नारायणाय नमः & नाभिं पूजयामि\\
९. & दामोदराय नमः & उदरं पूजयामि\\
१०. & माधवाय नमः & मध्यं पूजयामि\\
११. & वनमालिने नमः & वक्षः पूजयामि\\
१२. & चतुर्भुजाय नमः & बाहून् पूजयामि\\
१३. & कमलाधवाय नमः & कण्ठं पूजयामि\\
१४. & केशवाय नमः & कपोलौ पूजयामि\\
१५. & चन्द्रवदनाय नमः & चुबुकं पूजयामि\\
१६. & मन्दस्मिताय नमः & मुखं पूजयामि\\
१७. & पद्मनेत्राय नमः & नेत्रे पूजयामि\\
१८. & नासाहीनाय नमः & नासिकां पूजयामि\\
१९. & श्रुतिमध्याय नमः & श्रोत्रे पूजयामि\\
२०. & लक्ष्मीप्रियाय नमः & ललाटं पूजयामि\\
२१. & श्यामाङ्गाय नमः & शिरः पूजयामि\\
२२. & सर्वेश्वराय नमः & सर्वाण्यङ्गानि पूजयामि\\
\end{longtable}

\sect{उत्तराङ्ग-पूजा}
\begin{center}

\twolineshloka*
{धूपं गृहाण देवेश दशाङ्गेन सुवासितम्}
{धूतपाप दयाराशे दुरितोन्मुक्तविग्रह}
\textbf{\devAya{} धूपमाघ्रापयामि।}
\medskip

\twolineshloka*
{साज्यवर्तिसमायुक्तं पूज्यपाद पुरातन}
{दीपं गृहाण दयया ज्योतीरूप जगत्प्रभो}
\textbf{\devAya{} दीपं दर्शयामि।}
\medskip

\twolineshloka*
{नानाभक्ष्यसमायुक्तं नानाफलसमन्वितम्}
{नैवेद्यं गृह्यतां देव नीरजाक्ष नमोऽस्तु ते}
\textbf{\devAya{} नैवेद्यं निवेदयामि।}
\medskip

\twolineshloka*
{पूगीफलं सकर्पूरं नागवल्लीदलानि च}
{मुखवासाय गोविन्द गृह्यतां गरुडध्वज}
\textbf{\devAya{} ताम्बूलं समर्पयामि।}
\medskip

\twolineshloka*
{कर्पूरैः कल्पितं देव करुणारससागर}
{नीराजनं नीरजाक्ष गृह्यतां गुणसागर}
\textbf{\devAya{} कर्पूरनीराजनं समर्पयामि। कर्पूरनीराजनानन्तरम् आचमनीयं समर्पयामि।}
\medskip

\twolineshloka*
{प्रकृष्टपापनाशाय प्रकृष्टपदलब्धये}
{प्रदक्षिणं करोमि त्वां ज्योतीरूप जनार्दन}
\textbf{\devAya{} प्रदक्षिणान् समर्पयामि।}
\medskip

\twolineshloka*
{नमस्तेऽस्तु हृषीकेश नारायण नतावन}
{ज्योतीरूप जगन्नाथ जनार्दन नमो नमः}
\textbf{\devAya{} नमस्कारान् समर्पयामि।}
\medskip

\twolineshloka*
{मन्त्रपुष्पं गृहाणेदं महादेवादिवन्दित}
{मधुसूदन देवेश महापापविनाशन}
\textbf{\devAya{} मन्त्रपुष्पं समर्पयामि।}
\medskip

\twolineshloka*
{नानारत्नसमायुक्तं वज्रनालसमन्वितम्}
{मुक्ताकेसरसंयुक्तं स्वर्णपुष्पं ददाम्यहम्}
\textbf{\devAya{} स्वर्णपुष्पं समर्पयामि।}
\medskip

\twolineshloka*
{अखण्डदीपे देवेश तव पूजा कृता मया}
{सफला भवतु क्षिप्रं त्वत्प्रसादाज्जनार्दन}

\twolineshloka*
{ऋणपातकदौर्भाग्यव्याधीन् मे विनिवारय}
{पदं देहि परं देव प्रसीद पुरुषोत्तम}
\textbf{\devAya{} प्रार्थनाः समर्पयामि।}
\medskip

\end{center}

\dnsub{मासान्ते अर्घ्य-दीप-दान-विधिः}

अद्य पूर्वोक्त-एवङ्गुण-विशेषण-विशिष्टायाम् अस्याम् शुभतिथौ श्रीभूमिनीलासमेत-श्रीमहाविष्णु-प्रीत्यर्थं
श्रीभूमिनीलासमेत-श्रीमहाविष्णु-पूजान्ते क्षीरार्घ्यप्रदानं दीपदानं च करिष्ये॥

\twolineshloka*
{श्रीभूमिनीलासहित महाविष्णो दयानिधे}
{गृहाणेदं मया दत्तं ज्योतीरूप नमोऽस्तु ते}
\textbf{श्रीभूमिनीलासमेताय नमः इदमर्घ्यम् इदमर्घ्यम् इदमर्घ्यम्।}

अनेन दीपपूजनेन अर्घ्यप्रदानेन च श्रीमहाविष्णुः प्रीयताम्॥

इत्युक्त्वा ब्राह्मणं सम्पूज्य,

\twolineshloka*
{दीपं प्रज्वलितं साज्यम् अज्ञानविनिवारकम्}
{लक्ष्मीकरं ब्राह्मणाय विष्णुप्रीत्यै ददाम्यहम्}

इति दीपं दत्त्वा ब्राह्मणं भोजयित्वा स्वयं भुञ्जीत॥ विच्छेदे सति दीपस्य कृच्छ्रमेकं समाचरेत्।
सायं प्रातः प्रतिदिनं यथाशक्ति प्रपूजयेत्॥ मासे पूर्णे महादीपं यथाविभवविस्तरं सम्पूज्य च ततो
दद्यादाचार्याय सदक्षिणम्। दशदानं ततो दद्यात् भूरिदानं यथाबलम्। ब्राह्मणान् पूजयेत्पश्चात्
षड्रसैः ओदनैस्तथा॥

\twolineshloka*
{एवं यः कुरुते भक्त्या सर्वपापैः प्रमुच्यते}
{भुक्त्वा भोगान् स देहान्ते विष्णोस्सायुज्यमाप्नुयात्}

॥ इति श्रीव्रतरत्नाकरे अखण्डदीपपूजाकल्पः सम्पूर्णः ॥

\closesub

\closesection
